\clearpage
\section{Problem Analysis}

Compared to \textbf{assignment 3}, the main problem is no longer the local actor
graph, but the fact that the same logic must now work across \textbf{multiple
cluster nodes}.

The critical points are the following:
\begin{itemize}
    \item The \textbf{alarm state} must still have a \textbf{single owner}. If multiple control units changed state independently, the system would become inconsistent.
    \item \textbf{Sensors} and \textbf{keypads} must remain simple distributed actors. They generate commands or events, but they do not own the alarm state.
    \item \textbf{Timing constraints} from assignment 3 must remain valid even when messages cross node boundaries.
    \item A recreated control unit loses its previous in-memory state, so the system needs an explicit \textbf{safe recovery mode}.
\end{itemize}

This makes the clustered assignment a problem of \textbf{state ownership},
\textbf{discovery}, and \textbf{failure semantics}. The new distributed code must
therefore answer three questions clearly:
\begin{itemize}
    \item which actor owns the alarm state;
    \item how remote actors find that owner;
    \item what happens when that owner is recreated.
\end{itemize}

The solution adopted in the final project is to keep the \textbf{Control Unit}
as the only authority, expose it as a \textbf{cluster singleton}, and force it
to restart in \texttt{STARTUP\_RECOVERY} after recreation.
